\section{Conclusion}
\label{sec:conclusion}

The conventional knowledge graph was designed for a writer who no
longer writes it alone. We presented \sys, a store that inverts the
four defaults that assumption baked in — refuse at the gate,
bitemporal everything, partitioned trust with non-widening
composition, governance inside — and \census, a deterministic
lifecycle that measures each inversion against planted ground truth
and against a real writer's recorded trace. The measured story is
consistent: strictness is affordable when its cost falls on agents who
can read a refusal and retry; refusals are worth recording as signed
facts precisely because they are the events audit needs; composition
can be made safe without being made silent; and a store whose rules
are as historical as its data can replay not just what it knew but
what it required. An external sufficiency benchmark run over the
store's own exported evidence closes the loop from the outside:
every property-level governance question stays answerable exactly
when the evidence warrants it, with zero overclaim — and the run's
first pass improved the store it measured, sealing attribution into
the signed verdict.

Three boundaries are stated rather than hidden: no comparative
query-performance claims (the evaluator's ceiling is characterized,
not raced); labels are not access control (a floor refuses a query, it
does not hide rows); and denials replay as attestation checks, not
re-derivations, because refusal-by-rollback deliberately keeps the
verdict and discards the attempt. Future work runs in two directions:
the controlled agent-arm experiment (multiple tasks and models, with a
competency-question oracle) and lifting GS1–GS6 from design principles
to a portable contract — a governed-store specification other
substrates can claim and this benchmark can score, extending SARC's
compilation one layer further down for stores we have not built.
